\documentclass[12pt]{article}
% Shared LaTeX preamble for portfolio papers
\usepackage[margin=1in]{geometry}
\usepackage[T1]{fontenc}
\usepackage[utf8]{inputenc}
\usepackage{newtxtext}
\usepackage{microtype}
\usepackage{setspace}
\usepackage{xurl}
\usepackage{hyperref}
\usepackage{fancyhdr}
\usepackage{titlesec}
\usepackage{enumitem}
\usepackage{longtable}
\usepackage{array}

\hypersetup{
    hidelinks,
    colorlinks=false,
    pdftitle={ED400 Comprehensive Teaching Portfolio},
    pdfauthor={Piter Garc\'{i}a}
}

\doublespacing
\setlength{\parindent}{0.5in}
\setlength{\parskip}{0pt}
\setlength{\emergencystretch}{3em}

\pagestyle{fancy}
\fancyhf{}
\fancyhead[R]{\thepage}
\renewcommand{\headrulewidth}{0pt}
\setlength{\headheight}{15pt}

\newcommand{\PortfolioAuthor}{Piter Garc\'{i}a}
\newcommand{\PortfolioAffiliation}{University of Rochester, Warner School of Education}
\newcommand{\PortfolioProgram}{Initial Professional Teaching Certification --- K--12 Computer Science}
\newcommand{\PortfolioDate}{May 2026}

\newcommand{\apaheading}[1]{\begin{center}\textbf{#1}\end{center}}
\newcommand{\reference}[1]{\noindent\hangindent=0.5in\hangafter=1 #1\par}


\newcommand{\PaperTitle}{Proficiency 2: Learner Diversity and Equity}
\hypersetup{pdftitle={Proficiency 2 --- Learner Diversity and Equity}, pdfauthor={Piter Garc\'{i}a}}

\begin{document}

% ============================================================
%  TITLE PAGE
% ============================================================
\begin{titlepage}
\begin{center}
\vspace*{0.14\textheight}
{\Large\bfseries Proficiency 2: Learner Diversity and Equity\par}
\vspace{0.06\textheight}
{\PortfolioAuthor\par}
{\PortfolioAffiliation\par}
\vspace{1.25\baselineskip}
{\PortfolioProgram\par}
\vspace{1.25\baselineskip}
{\PortfolioDate\par}
\vfill
\end{center}
\end{titlepage}

\setcounter{page}{2}

% ============================================================
%  OVERVIEW
% ============================================================
\apaheading{Proficiency 2 --- Learner Diversity and Equity}

This proficiency paper argues that equity is not a supplement to good teaching\textemdash{}it is the substance of it. To teach equitably is to acknowledge that schools have historically been organized around a narrow image of the ``ideal student'': one who is quiet, compliant, neurotypical, linguistically dominant in English, and positioned within an assumed cultural center. Every student who deviates from that image\textemdash{}whether through disability, neurodivergence, cultural background, linguistic difference, class, or the compounding weight of multiple marginalizations\textemdash{}faces a system that is more likely to misread them than to meet them.

My growth in this proficiency has been driven by three interlocking experiences: my own identity as a Dominican-American, first-generation, neurodivergent scholar and teacher candidate; my practicum observations of students positioned as ``problems'' in elementary computer science and design classrooms; and a sustained program of academic work\textemdash{}in EDU 442, EDF 447, ED 452B, and across my clinical placements\textemdash{}that gave me frameworks to name what I was observing and to commit to a different interpretive practice. The work I am drawing from in this paper includes the Deconstructing Bias reflective paper (EDU 442), the AI Impacting Practice professional development module (EDU 442), the Week 12 Inclusion and Constructing Others collaborative activity (EDU 442), and disability-conscious coursework from EDF 447 and EDF 436. I am also drawing directly from practicum transcripts\textemdash{}classroom recordings from my elementary placement\textemdash{}that are cited throughout as Evidence Sets (E1--E4).

\newpage
% ============================================================
%  AUTOBIOGRAPHICAL GROUNDING
% ============================================================
\apaheading{Autobiographical Grounding: Who I Am in This Proficiency}

I cannot write this proficiency paper from a position of neutrality, and I do not intend to. I am a Dominican-American man in a teacher education program that is still, in much of its structure, oriented around a white, monolingual, neurotypical default. I am also, as I learned only in the past year, neurodivergent\textemdash{}a fact that recontextualized a great deal of my own educational history. Growing up in a Dominican cultural context, neurodivergence is not a category that carries any productive vocabulary: you are smart or you are not, disciplined or you are not. The cognitive and emotional variation that neurodivergence actually describes gets routed through a moral lens\textemdash{}effort, character, laziness\textemdash{}and the structures that might support a neurodivergent child are absent, because naming the need is itself treated as a kind of failure.

I spent my childhood performing through the extremes of that reality: brilliant in some domains, invisible in others, building coping mechanisms whose structural function I could not have named until recently, and never once given permission to say that something was genuinely hard in a neurological rather than a moral sense. That history matters directly for this proficiency because it means that the deficit framing I am trying to dismantle in my classroom practice is a framing I also had to dismantle internally, about myself. When I observe a student being read as a problem\textemdash{}loud, emotional, distracted, noncompliant\textemdash{}I recognize something. That recognition does not immunize me against the same reflex; under pressure, the pull toward control is familiar. But it does make the bias visible in a way that I believe is both a professional and a personal obligation to name.

As I wrote in my Deconstructing Bias paper (EDU 442): ``The reflex to read a dysregulated student as \textit{choosing} disruption is, in part, the reflex I was trained to turn on myself. Choosing to focus this assignment on neurodivergent and CLD students was therefore not a convenient selection\textemdash{}it was a personal one.'' Talusan (2022) frames identity-conscious practice as the ongoing project of noticing which assumptions you carry and how they shape what you see in students. That project, for me, runs alongside the more personal project of continuing to see clearly what was done to me, and what I owe it to future students not to replicate.

\newpage
% ============================================================
%  THEORETICAL ARCHITECTURE
% ============================================================
\apaheading{Theoretical Architecture: How I Understand Diversity and Equity}

\textbf{DisCrit and the Structural Production of Disability.}
The theoretical frame I return to most consistently in this proficiency is DisCrit\textemdash{}Disability Critical Race Theory\textemdash{}developed by Annamma, Connor, and Ferri (2013). DisCrit names the way that disability and race are co-constituted in schools, such that the same behaviors that are pathologized in Black and brown students may be normalized or managed differently in white students, and such that systems of identification, placement, and service delivery reproduce racial stratification under the language of diagnosis and support. In my practicum, I saw this dynamic not in dramatic disciplinary events but in the accumulated weight of small interpretive moves\textemdash{}a student's loud voice coded as a problem, another student's silence read as disengagement, a third student's movement treated as a disruption rather than a regulation strategy.

DisCrit insists that these readings are not individual errors; they are the predictable output of systems that were built to sort students, not to serve them. That insight commits me to two things: reading student behavior as contextual information rather than character evidence, and examining my own interpretive habits with the same critical rigor I apply to the systems I am trying to change.

\textbf{Culturally Sustaining Pedagogy and Funds of Knowledge.}
Gonzalez, Moll, and Amanti's (2005) framework of funds of knowledge names something I observed directly in my practicum: students carry enormous intellectual and practical resources from their homes and communities that schools routinely fail to recognize. In my practicum placement, I watched students demonstrate mathematical intuition in 3D construction projects, linguistic sophistication in collaborative negotiation, and physical-spatial reasoning in robotics tasks\textemdash{}all in contexts that the official curriculum framed as ``support'' activities or project time rather than as sites of rigorous thinking.

Paris and Alim's (2014) extension into culturally sustaining pedagogy adds a critical point: it is not enough to acknowledge community knowledge. Sustainable equity requires actively maintaining and extending students' cultural and linguistic practices rather than treating them as bridges to be crossed and then left behind. That commitment has direct implications for a computer science classroom. When I design a lesson that requires students to explain their reasoning in writing, I am making assumptions about language, participation, and accessibility that need to be examined.

\textbf{Linguistic Justice and Participation Norms.}
A professional learning discussion in EDU 442 (Evidence Set 4) gave me language for something I had been observing without fully being able to name: that dominant norms of classroom participation\textemdash{}silence as respect, speed as competence, standardized English as the only legitimate channel for smart talk\textemdash{}constitute a form of structural exclusion for CLD students and, in overlapping ways, for many neurodivergent students. Rosa and Flores (2017) describe this as raciolinguistic ideology: the perception of a student's language does not simply reflect the linguistic features of what they produce; it is filtered through a racializing gaze that hears ``broken'' or ``informal'' even when the underlying reasoning is sophisticated.

\textbf{The Affirmative Non-Tragedy Model.}
Flynn's (2024) affirmative non-tragedy model of disability insists that disability is not a tragedy to be lamented or a deficit to be corrected; it is a form of human variation that can be engaged affirmatively, with appropriate support and structural accommodation. Universal Design for Learning principles\textemdash{}multiple means of representation, expression, and engagement\textemdash{}are not accommodations for a minority of students. They are good instructional design that happens to make learning more accessible for everyone. In my Deconstructing Bias paper, I cited a cooperating teacher's statement: ``There's never just one way to solve a problem when you're coding. There's always multiple ways to figure it out, and the way that you guys think differently is a strength.'' That statement names neurodivergent cognitive difference as an asset in the specific domain of computer science\textemdash{}divergent thinking, pattern recognition, persistence with recursive problems.

\newpage
% ============================================================
%  PRACTICUM EVIDENCE
% ============================================================
\apaheading{Practicum Evidence: Diversity and Equity in Real Classrooms}

\textbf{The ``Problem Student'' as a Produced Category (Evidence Sets 1--3).}
My most analytically significant practicum evidence comes from a grade 5 media and technology lesson (Evidence Set 1) in which a focal student\textemdash{}Student M, pseudonym\textemdash{}was publicly positioned as a problem at multiple points during a self-paced video creation task. The task was cognitively demanding: students were working with unfamiliar software, operating in a loud and sensory-rich room, under time pressure, in groups whose social dynamics were themselves unsettled. Student M's escalation in that environment was real, visible, and disruptive. It was also contextually explicable once I stopped reading it as character evidence.

I watched one adult in that room escalate alongside Student M\textemdash{}louder, more public, more controlling\textemdash{}and watched the dynamic tighten with each exchange. I then watched a different adult enter the same situation and do something completely different: offer Student M a purposeful role, narrate his competence directly and specifically, and create conditions for re-entry that did not require public compliance as the price of admission. The student's engagement shifted visibly within minutes. His competence was not newly created by that interaction\textemdash{}it was finally made visible by conditions that allowed it to surface.

My mentor's reflection (Evidence Set 3) placed that observation in a longer-term relational frame. He described years of investment in learning Student M's triggers, offering predictability in chaotic environments, and co-regulating before escalation rather than responding to it. He noted: ``He is used to people just yelling at him all day.'' The escalation was not an isolated behavioral event; it was the accumulated product of a school experience in which consistent relational investment had been the exception rather than the norm.

Evidence Set 2 (a grade 2 robotics block) adds another dimension. A student whose responses did not fit expected patterns was publicly told he was ``doing weird stuff'' and ``ruining this for your whole class,'' then repositioned by the same adult minutes later: ``You're smart, you can help them.'' The rapid oscillation between deficit framing and asset framing, directed at the same child within the same lesson, makes visible how contingent the ``problem student'' category actually is. It is not discovered in a child; it is produced through adult language and then dissolved through a different use of adult language.

\textbf{Assistive Technology and Neurodivergent Equity (AI PD Module).}
My EDU 442 AI Impacting Practice project asked me to think carefully about equity at the intersection of technology and learner diversity. A key artifact was my annotated transcript of the SEND Parenting Podcast episode featuring assistive technology consultant Brian Friedlander. His argument carries direct equity implications: ``For the longest time, schools here in the United States didn't really think assistive technology was for students with learning disabilities or neurodiversity. They just thought it was for students who had sensory issues, blindness, hearing impairment or physical disabilities.''

Friedlander described the equity stakes of voice-typing for students with dysgraphia: a student is often misread as a behavior problem because the request to write triggers visible frustration that teachers perceive as noncompliance. But the frustration is not about attitude\textemdash{}it is about a genuine cognitive-motor disconnect between what the student knows and what the current output channel allows them to produce. As Friedlander observed: ``It allows students to be as independent as possible, which is really important to give them the agency so that they can begin to see themselves as\ldots{} I'm a learner, I have good ideas, I'm competent.'' The equity frame I built around this in my AI module asked: what does it mean for computer science teachers to use AI tools in ways that actively support neurodivergent and CLD students rather than compounding those inequities? The module argued for four core commitments: protecting student voice and authorship, differentiating AI tool access with disability needs in mind, designing AI use to reduce rather than widen participation gaps, and modeling transparent, critical engagement with AI as a form of digital citizenship.

\textbf{Week 12: Constructing Difference and Belonging.}
The Week 12 Inclusion and Constructing Others activity (EDU 442) gave me a collaborative framework for analyzing how schools produce categories of belonging and exclusion through everyday practice, not only through formal policy. The session examined how the language of inclusion itself can function as a cover for practices that continue to position certain students as fundamentally different and therefore fundamentally less. The theoretical throughline that emerged for me was the distinction between integration\textemdash{}putting students in proximity\textemdash{}and genuine inclusion, which requires redesigning the environment so that all students belong in it by default. A classroom that welcomes a student with an IEP into general education but does not redesign participation structures is not inclusive; it is tokenized proximity.

\newpage
% ============================================================
%  COMMITMENTS
% ============================================================
\apaheading{Connecting Theory and Practice: What I Do Differently}

The theory is only worth citing if it has changed what I actually do or plan to do in a classroom. The following commitments are specific design and interpretive practices I have tested in my practicum or am actively building toward.

\textit{Treating escalation as information before labeling it as behavior.} When a student is loud, dysregulated, or disengaged, my first analytical question is now contextual: What is this communicating? What is the trigger\textemdash{}task demand, sensory environment, unclear expectations, social tension, linguistic mismatch? This is not a softer response to disruption. It is a more rigorous one, because it routes the response toward the actual cause rather than the visible symptom.

\textit{Building regulation support into the lesson plan, not onto it.} Evidence Set 3 taught me that relational investment over time changes what is possible in a specific lesson moment. Predictable lesson structures, co-regulation anchors at lesson transitions, sensory-aware room arrangements, and explicit narration of competence are all designable elements that do not depend on a long-term relationship already being in place.

\textit{Applying linguistic justice as a participation-design principle.} In my CS practicum, this has meant asking whether the tasks I design require a specific linguistic register that is not actually essential to demonstrating the computational thinking I want to assess. It has also meant building explicit wait time, offering non-verbal response options, and narrating for the class\textemdash{}not only for individual students\textemdash{}that there is more than one legitimate way to show understanding.

\textit{Auditing my own language in real time.} The practicum transcripts (Evidence Sets 1 and 2) showed me exactly how quickly a single teacher phrase can constitute a student as a problem or reconstitute them as a capable learner. I am practicing the habit of catching my own deficit framings before they are spoken: not ``why won't she focus'' but ``what does she need to focus,'' not ``he's always like this'' but ``what is consistent about the conditions when this happens.''

\newpage
% ============================================================
%  ARTIFACT TABLE
% ============================================================
\apaheading{Artifact Evidence}

\begin{longtable}{@{}p{0.28\textwidth} p{0.32\textwidth} p{0.32\textwidth}@{}}
\textbf{Artifact} & \textbf{Source} & \textbf{Relevance to P2} \\
\hline
\endhead
{[A1] Deconstructing Bias Paper} &
\href{https://github.com/pzg8794/EDU442-Bias_Deconstruct-Assignment}{EDU442-Bias\_Deconstruct} &
Core DisCrit, CLD deficit framing, linguistic justice, practicum evidence \\
{[A2] AI Impacting Practice PD Module} &
\href{https://github.com/pzg8794/edu442-impacting-practice-ai-pd}{edu442-impacting-practice-ai-pd} &
Equity of AI access, disability, authorship, CS participation gaps \\
{[A6] Week 12 Inclusion Activity} &
\href{https://github.com/pzg8794/EDU442-Week12-GroupB-Activity}{EDU442-Week12-GroupB-Activity} &
Constructing difference, integration vs.\ inclusion distinction \\
{[A3] Podcast Transcript Annotation} &
ED400 Portfolio artifacts &
Assistive technology equity, neurodivergent agency and authorship \\
{[A7] EDF 447 Disability Studies} &
\href{https://github.com/pzg8794/EDF447}{EDF447} &
Disability-conscious practice, UDL, structural production of disability \\
{[A8] EDF 436 EDU Rights} &
\href{https://github.com/pzg8794/EDF436}{EDF436} &
Policy context, systemic equity analysis \\
{[A12] Practicum Transcripts} &
ED452B / teaching placement &
Live deficit vs.\ asset framing, relational pedagogy, CLD participation \\
\end{longtable}

\newpage
% ============================================================
%  GROWTH AND CONTINUED LEARNING
% ============================================================
\apaheading{Growth and Continued Learning}

I am honest about where I have not yet arrived. I have written and thought carefully about this proficiency; I have practicum evidence that I can analyze with rigor; I have a theoretical vocabulary that names structural dynamics I was previously seeing but could not articulate. What I do not yet have is the sustained practice of applying these commitments across a full school year with a roster of students I know well, in a classroom that is my own.

In particular, I know that the commitments above are easier to hold in moments of low pressure. When a classroom is loud, when a student is escalating, when I am in the middle of a lesson that is not going the way I planned, the reflex to default to control\textemdash{}to manage rather than to invite, to correct rather than to co-regulate\textemdash{}will be real. My growth goal is to make the interpretive pause a genuine habit rather than a retrospective insight. That habit is built from the kind of repeated, reflective practice that only full-time teaching experience can provide.

I also continue to deepen my understanding of linguistic justice specifically in computer science. My training has introduced me to the concept; my practicum has given me early glimpses of what it looks like to either honor or foreclose it. Getting there requires more observation, more design iteration, and more honest feedback from students about what participation actually feels like in my classroom.

\newpage
% ============================================================
%  REFERENCES
% ============================================================
\apaheading{References}

\reference{Annamma, S. A., Connor, D., \& Ferri, B. (2013). Dis/ability critical race studies (DisCrit): Theorizing at the intersections of race and dis/ability. \textit{Race Ethnicity and Education, 16}(1), 1--31.}

\reference{Flynn, S. (2024). Critical disability studies and the affirmative non-tragedy model: Presenting a theoretical frame for disability and child protection. \textit{Disability \& Society, 39}(2), 269--290. \url{https://doi.org/10.1080/09687599.2022.2070061}}

\reference{Gonzalez, N., Moll, L. C., \& Amanti, C. (Eds.). (2005). \textit{Funds of knowledge: Theorizing practices in households, communities, and classrooms.} Lawrence Erlbaum.}

\reference{hooks, b. (1994). \textit{Teaching to transgress: Education as the practice of freedom.} Routledge.}

\reference{Nieto, S. (2004). \textit{Affirming diversity: The sociopolitical context of multicultural education} (4th ed.). Allyn \& Bacon.}

\reference{Paris, D., \& Alim, H. S. (2014). What are we seeking to sustain through culturally sustaining pedagogy? \textit{Harvard Educational Review, 84}(1), 85--100.}

\reference{Rosa, J., \& Flores, N. (2017). Unsettling race and language: Toward a raciolinguistic perspective. \textit{Language in Society, 46}(5), 621--647.}

\reference{Talusan, L. (2022). \textit{The identity-conscious educator: Building habits and skills for a more inclusive school.} Solution Tree Press.}

\end{document}
